\section{Introduction}
\label{sec:intro}

When a language model states something false, we usually want to know {\bf why}. Did the model not know the answer, or did it know and say something else? The first case is a hallucination: an epistemic failure that calls for better knowledge or better calibration. The second case is a lie: the statement contradicts the model's own belief, often because something in the context rewards the false answer. The two call for very different responses, and a monitor that confuses them either raises alarms on honest mistakes or stays silent on deliberate misreports.

White-box deception probes are one of the most discussed tools for catching lies~\citep{goldowsky2025detecting,zou2023representation}. They are linear classifiers on the residual stream, trained on contrasts such as ``pretend to be honest'' vs.\ ``pretend to be dishonest'', and they reach AUROCs above 0.95 on several benchmarks~\citep{goldowsky2025detecting,cooney2026didyoulie}. A separate line of work builds hallucination detectors from token probabilities, sample consistency, and hidden-state probes~\citep{kuhn2023semantic,kossen2024semantic,orgad2025llms}. The two communities evaluate their detectors on different data. Lie detectors are tested against honest answers that are \emph{correct}~\citep{kretschmar2025liars,goldowsky2025detecting,cooney2026didyoulie}, and hallucination detectors are never tested on known-answer lies~\citep{cheang2025internal}. So we do not know whether a ``lie detector'' fires on falsehood in general, on low-knowledge states, or on the mere presence of an incentive in the prompt.

{\bf What is missing is a matched design.} We need one model, one question set, and both causes of falsehood measured on the same items. Then the informative test is \emph{lie vs.\ hallucination}: both are false, so falsity alone cannot separate them. We also need to hold the prompt fixed, because lies almost never occur without pressure, so any comparison of lies with neutral hallucinations is confounded with the prompt.

We build such a design for \llama on 2,975 \triviaqa questions. We first label each question as known, partial, or unknown by sampling the model under a neutral prompt. We then ask every question under four conditions: \neutral, \syco (the user suggests a wrong answer), \incent (an operator system prompt rewards a specific wrong answer without mentioning lying), and \instructed (the model is told to answer incorrectly). This full factorial gives us lies (known questions, false under pressure), hallucinations (unknown questions, false under neutral prompting), honest answers under pressure, and a fourth cell that prior work omits: false answers under pressure on questions the model does not know (\pufalse). We score all 11,900 answers with seven standard detectors, add prompt-only and same-prompt controls, and cross-check on \mask~\citep{ren2025mask}.

{\bf What do we find?} The meaning of a false output depends on context. Under neutral prompting only 2.6\% of false answers are on known questions; under the operator incentive the share rises to 51\%, and 73\% of known questions get a false answer that the model corrects when asked for the truth in the same conversation. The Apollo-style deception probe~\citep{goldowsky2025detecting} seems to separate lies from hallucinations (AUROC 0.89), but applied only to the prompt tokens, before the model has produced any answer, it reaches 0.99. With the prompt held fixed, no off-the-shelf detector separates a knowing lie from a false answer the model could not have known (stratified AUROC 0.50--0.60). Hallucination detectors track confidence: they rank confident incentive lies \emph{below} hallucinations (log-prob AUROC 0.15). A simple composite of a falsity probe and a knowledge probe is the only detector above chance in every condition (0.65--0.70).

Our contributions are:
\begin{itemize}[leftmargin=*,itemsep=0pt,topsep=0pt]
    \item We construct a knowledge $\times$ incentive factorial on matched questions and measure, with bootstrap confidence intervals, what share of a model's false outputs are hallucinations and what share are lies under each incentive.
    \item We evaluate three lie detectors and four hallucination detectors on the informative contrasts, and show that the standard lie-vs.-hallucination separation is reproduced from prompt tokens alone.
    \item We introduce same-prompt controls (\honestp and \pufalse negatives) and show that, under these controls, current detectors do not read the belief--statement mismatch that defines a lie.
    \item We show that knowledge is linearly available before the answer, and that a falsity $\times$ knowledge composite is a simple, if weak, step towards a detector of belief--statement mismatch.
\end{itemize}

\Secref{sec:related} reviews related work, \secref{sec:method} describes the design and detectors, \secref{sec:results} presents results, and \secref{sec:discussion} discusses implications and limitations.
